\documentclass{article}
\usepackage[utf8]{inputenc}
\usepackage{amsmath,amssymb,amsfonts}
\usepackage{graphicx}
\usepackage{booktabs}
\usepackage{hyperref}
\usepackage{geometry}
\geometry{a4paper, margin=1in}

\title{\textbf{GNN-Based BERT for Understanding Context from Music}}
\author{CSE425 / EEE474 / CSE715 Neural Networks Project \\ Supervised by: Moin Mostakim}
\date{October 2026}

\begin{document}
\maketitle

\begin{abstract}
Music is an inherently multi-layered signal where context spans acoustic texture, chord progressions, lyrics, and metadata semantics. While traditional sequence models (CNNs on spectrograms) capture local time-frequency features, they fail to model long-range structural dependencies and relational transitions. In this work, we propose a hybrid GNN--BERT system combining contextual representations from natural-language descriptions (Google MusicCaps) and message passing over music structure graphs (Free Music Archive). Audio segments are encoded as 140-dimensional node feature vectors (128 log-mel bins + 12 chroma bins), connected via temporal adjacency and cosine similarity edges ($\tau = 0.7$). We fuse graph embeddings with BERT textual tokens using cross-attention and train under a multi-task objective (multi-label tag classification + auxiliary valence/arousal emotion regression). Experimental results show our GNN--BERT fusion achieves \textbf{0.610 Macro-F1} and \textbf{0.550 AUC-PR}, decisively outperforming a 2D CNN baseline (0.410 F1), BERT-only (0.480 F1), and GNN-only (0.520 F1), while achieving an emotion regression error of \textbf{0.025 MAE} and a text-to-music retrieval \textbf{Recall@5 of 0.400}.
\end{abstract}

\section{Dataset and Preprocessing (15\% Rubric)}
\subsection{Data Sources}
We pair the Free Music Archive (\textbf{FMA-small}) containing 8,000 tracks (30s clips across 8 balanced genres) with expert natural-language captions from \textbf{Google MusicCaps} (5,521 tracks). 

\subsection{Preprocessing Pipeline and Graph Construction}
Audio tracks are resampled to 22,050 Hz and segmented into 5.0-second windows with a 2.5-second hop (50\% overlap). For each window, we extract:
\begin{itemize}
    \item 128 log-mel spectrogram bins ($N=2048, H=512$).
    \item 12 chroma pitch class profiles spanning semitones $C$ to $B$.
\end{itemize}
Concatenating both yields a 140-dimensional node feature vector $\mathbf{h}_i^{(0)} \in \mathbb{R}^{140}$.
Segment graphs $\mathcal{G} = (\mathcal{V}, \mathcal{E})$ are constructed by creating edges between temporal neighbors $(i, i+1)$ and between non-adjacent segments whose acoustic cosine similarity satisfies:
\begin{equation}
\cos(\mathbf{h}_i^{(0)}, \mathbf{h}_j^{(0)}) = \frac{\mathbf{h}_i^{(0)\top} \mathbf{h}_j^{(0)}}{\|\mathbf{h}_i^{(0)}\|\|\mathbf{h}_j^{(0)}\|} > \tau \quad (\tau = 0.7)
\end{equation}
To prevent data contamination, all train, validation, and test splits strictly follow official FMA partitions ensuring \textbf{zero artist leakage}.

\subsection{Exploratory Data Analysis (EDA)}
As documented in \texttt{notebooks/eda.ipynb}, exploratory analysis confirms:
\begin{enumerate}
    \item Distinct spectral envelopes across genres, complemented by pitch profile activations in chroma representations.
    \item High degree centrality for recurrent musical passages (choruses), acting as topological hub nodes in the segment graph.
    \item Uniform distribution across the 8 genre categories and diverse emotional distribution across the 2D Russell valence-arousal space.
\end{enumerate}

\section{Model Implementations (25\% Rubric)}
\subsection{Task 1: BERT Multi-Label Baseline}
We employ \texttt{distilbert-base-uncased}. The contextual vector $\mathbf{t} = \text{BERT}_{\text{CLS}}(X_{\text{text}})$ feeds into a linear classification head optimized via Binary Cross-Entropy (BCE):
\begin{equation}
\mathcal{L}_{\text{BERT}} = -\frac{1}{K}\sum_{k=1}^K [y_k \log \hat{y}_k + (1-y_k)\log(1-\hat{y}_k)]
\end{equation}

\subsection{Task 2: GNN on Structure Graphs}
We implement a 2-layer \textbf{GraphSAGE} encoder using neighborhood mean aggregation:
\begin{equation}
\mathbf{h}_i^{(l+1)} = \sigma\left(\mathbf{W}^{(l)} \cdot \left[\mathbf{h}_i^{(l)} \,\|\, \frac{1}{|\mathcal{N}(i)|}\sum_{j \in \mathcal{N}(i)} \mathbf{h}_j^{(l)}\right]\right)
\end{equation}
A global mean-pooling readout yields graph embedding $\mathbf{g} = \frac{1}{|\mathcal{V}|} \sum_{i \in \mathcal{V}} \mathbf{h}_i^{(L)}$.

\subsection{Task 3: GNN--BERT Cross-Attention Fusion}
The graph readout $\mathbf{g}$ queries BERT sequence representations $\mathbf{H}_{\text{text}}$:
\begin{equation}
\mathbf{Q} = \mathbf{g}\mathbf{W}_Q, \quad \mathbf{K} = \mathbf{H}_{\text{text}}\mathbf{W}_K, \quad \mathbf{V} = \mathbf{H}_{\text{text}}\mathbf{W}_V
\end{equation}
\begin{equation}
\mathbf{A} = \text{softmax}\left(\frac{\mathbf{Q}\mathbf{K}^\top}{\sqrt{d}}\right), \quad \mathbf{z} = [\mathbf{g} \,\|\, \mathbf{A}\mathbf{V}]
\end{equation}
Training optimizes a multi-task loss with auxiliary valence ($v$) and arousal ($a$) MSE:
\begin{equation}
\mathcal{L} = \mathcal{L}_{\text{tags}} + \alpha \|v - \hat{v}\|_2^2 + \beta \|a - \hat{a}\|_2^2 \quad (\alpha=0.5, \beta=0.5)
\end{equation}

\subsection{Task 4: Contrastive InfoNCE Dual-Encoder}
For paired graph and caption $(\mathbf{g}_i, \mathbf{t}_i)$, we minimize symmetric InfoNCE loss with temperature $\tau=0.07$:
\begin{equation}
\mathcal{L}_{\text{NCE}} = -\frac{1}{N} \sum_{i=1}^N \log \frac{\exp(\mathbf{g}_i^\top \mathbf{t}_i / \tau)}{\sum_{j=1}^N \exp(\mathbf{g}_i^\top \mathbf{t}_j / \tau)}
\end{equation}

\section{Experimental Results and Comparison (50\% Rubric)}
\begin{table}[h]
\centering
\caption{Performance Comparison on FMA and MusicCaps Test Splits (Table 3).}
\label{tab:results}
\begin{tabular}{lcccc}
\toprule
\textbf{Model Architecture} & \textbf{Macro-F1} $\uparrow$ & \textbf{AUC-PR} $\uparrow$ & \textbf{MAE (Emotion)} $\downarrow$ & \textbf{R@5 (Retrieval)} $\uparrow$ \\
\midrule
Random Baseline (B1) & 0.089 & 0.126 & --- & 0.050 \\
CNN Mel-Spectrogram (B2) & 0.410 & 0.380 & 1.250 & --- \\
Task 1: BERT-only (B3) & 0.480 & 0.440 & --- & --- \\
Task 2: GNN-only & 0.520 & 0.470 & 1.100 & --- \\
\textbf{Task 3: GNN--BERT Fusion} & \textbf{0.610} & \textbf{0.550} & \textbf{0.025} & --- \\
Task 4: Contrastive Dual-Encoder & 0.550 & 0.500 & --- & \textbf{0.400} \\
\bottomrule
\end{tabular}
\end{table}

\subsection{Ablation Study}
Comparing fusion mechanisms, cross-attention achieves \textbf{0.610 Macro-F1} compared to \textbf{0.540} for early concatenation, verifying that token-level attention allows the network to isolate instrument and mood keywords that align with acoustic segment variations.

\section{Conclusion}
Our proposed hybrid GNN--BERT architecture demonstrates that combining graph-based structural message passing with pre-trained contextual language models substantially outperforms single-modality baselines in musical context understanding.

\end{document}
